Claims about activation interventions are claims about specific
direction vectors, and about the random directions used as their
controls. We test both. First, we give a small provenance protocol that
stores each direction as a hashed, hermetically copied artifact, and
evaluate it by injecting 13 realistic faults into a real experiment. A
hardened self-contained integrity gate catches 12 (the original caught
6). The 13th, a run silently redone end to end, is caught only by a hash
ledger pinned outside the run. A common regenerate-then-verify CI step
catches 2. Second, in a pre-registered study on GPT-2, Qwen3-0.6B and
Llama-3.2-1B (three blocks each, 16 random directions per control and
condition), random directions are not reliably null controls. With a
heterogeneity test we calibrated on the real data (the pre-registered $\chi^2$
test proved anticonservative, rejecting true nulls 17--23\% of the
time), random edits matched to a fixed general-text KL budget are
heterogeneous in 23 of 27 cells, and additive edits in 27 of 36. Plain
unit projections are null at the first block of every model but
heterogeneous deeper in (14 of 27 cells overall). The learned bias
direction itself beats all 48 of its random controls in only 9 of 30
condition cells. Support-avoiding and subspace-matched random directions
behave like unconstrained ones, and Jacobian-lens propagation does not
predict which directions matter. Finally, regenerating directions
reproduced every hash on the same machine seven months later, but not on
different hardware, where vectors shifted by a few float32 ULPs while
results agreed to $10^{-6}$. Stored vectors, not regenerated ones, are what
should be pinned.
